% ============================================================
\section{Day 2}
% ============================================================

\subsection{Readings}

\begin{itemize}
    \item Bernhard and Michael Grzimek, \textit{Serengeti Shall Not Die}.
\end{itemize}



% ============================================================
\subsection{Theory of the mean}
% ============================================================

A useful way of constructing a dynamical model is to keep track of the mean value of one or more variables of interest.

The basic protocol is:

\begin{enumerate}
    \item Identify the variables of interest.
    \item Enumerate all the ways in which those variables can increase or decrease, and assign rates to those processes.
    \item Write an update rule describing the change during a small time interval $\Delta t$.
    \item Take the continuum limit and integrate the resulting differential equation to obtain the time evolution.
\end{enumerate}


% ------------------------------------------------------------
\subsubsection{Example: ligand-receptor binding}
% ------------------------------------------------------------

Consider a ligand $L$ that can bind reversibly to a receptor $R$ to form a ligand--receptor complex $LR$.

The variables of interest are therefore the concentrations

\[
    [R], \qquad [L], \qquad [LR].
\]

Here $[R]$ denotes the concentration of free receptors, $[L]$ the concentration of free ligands, and $[LR]$ the concentration of bound ligand--receptor complexes.
% 
% 
% ------------------------------------------------------------
\medskip
\\
\textbf{Step 1: identify the variables}
% ------------------------------------------------------------
\\
A schematic representation of the three states is
\begin{figure}[h]
    \centering
    \begin{tikzpicture}[
        ligand/.style={circle, draw, minimum size=0.7cm, inner sep=0pt},
        receptor/.style={rounded corners=2pt, draw, minimum width=1.0cm, minimum height=0.7cm},
        reaction/.style={Latex-Latex, thick}
    ]
        \node[receptor] (R) at (0,0) {$R$};
        \node at (1.0,0) {$+$};
        \node[ligand] (L) at (2.0,0) {$L$};

        \draw[reaction] (2.8,0) -- (5.2,0);
        \node at (4.0,0.28) {$k_{\mathrm{on}}$};
        \node at (4.0,-0.28) {$k_{\mathrm{off}}$};

        \node[receptor] (Rc) at (6.2,0) {$R$};
        \node[ligand] (Lc) at (6.2,0.68) {$L$};
        \node at (7.1,0) {$LR$};
    \end{tikzpicture}
    \caption{Reversible binding of a ligand $L$ to a receptor $R$ to form the complex $LR$.
    The same ligand and receptor symbols are used in the unbound and bound states.}
\end{figure}
% 
% 
% ------------------------------------------------------------
\medskip
\\
\textbf{Step 2: enumerate the processes and their rates}
% ------------------------------------------------------------
\\
The reaction is
\[
    R + L \xrightleftharpoons[k_{\mathrm{off}}]{k_{\mathrm{on}}} LR.
\]
The forward reaction is ligand binding.
According to the law of mass action, its rate is proportional to both the free ligand concentration and the free receptor concentration,
\[
    \text{binding rate} = k_{\mathrm{on}}[L][R].
\]

The reverse reaction is ligand unbinding.
Its rate is proportional to the concentration of bound complexes,
\[
    \text{unbinding rate} = k_{\mathrm{off}}[LR].
\]
The association rate constant $k_{\mathrm{on}}$ has units of concentration$^{-1}$ time$^{-1}$.
The dissociation rate constant $k_{\mathrm{off}}$ has units of time$^{-1}$.
% 
% 
% ------------------------------------------------------------
\medskip
\\
\textbf{Step 3: write the update rule}
% ------------------------------------------------------------
\\
During a short time interval $\Delta t$, the concentration of bound complexes changes according to
\[
    [LR](t+\Delta t) = [LR](t) + k_{\mathrm{on}}\Delta t[L][R] - k_{\mathrm{off}}\Delta t[LR].
\]
The first additional term represents newly formed complexes.
The second represents complexes lost through dissociation.
Subtracting $[LR](t)$ from both sides and dividing by $\Delta t$ gives
\[
    \frac{[LR](t+\Delta t)-[LR](t)}{\Delta t} = k_{\mathrm{on}}[L][R] - k_{\mathrm{off}}[LR].
\]
Taking the limit $\Delta t\rightarrow0$ gives
\[
    \boxed{\frac{\dd[LR]}{\dd t} = k_{\mathrm{on}}[L][R] - k_{\mathrm{off}}[LR].}
\]
% 
% 
% ------------------------------------------------------------
\medskip
\\
\textbf{Steady state}
% ------------------------------------------------------------
\\
At steady state, the mean concentration of bound complexes is no longer changing,
\[
    \frac{\dd[LR]}{\dd t} = 0.
\]
Therefore,
\[
    k_{\mathrm{on}}[L][R] - k_{\mathrm{off}}[LR] = 0.
\]
Hence
\[
    k_{\mathrm{on}}[L][R] = k_{\mathrm{off}}[LR].
\]
Rearranging,
\[
    \frac{[L][R]}{[LR]} = \frac{k_{\mathrm{off}}}{k_{\mathrm{on}}}.
\]
We define the dissociation constant
\[
    \boxed{K_d = \frac{k_{\mathrm{off}}}{k_{\mathrm{on}}}.}
\]
Therefore,
\[
    K_d = \frac{[L][R]}{[LR]}.
\]
Equivalently,
\[
    \frac{[LR]}{[R]} = \frac{[L]}{K_d}.
\]
% 
% 
% ------------------------------------------------------------
\medskip
\\
\textbf{Probability that a receptor is bound}
% ------------------------------------------------------------
\\
The total receptor concentration is the sum of free and bound receptors,
\[
    [R]_{\mathrm{tot}} = [R] + [LR].
\]
The fraction of receptors that are bound is therefore
\[
    P_{\mathrm{bound}} = \frac{[LR]}{[R]+[LR]}.
\]
Using
\[
    \frac{[LR]}{[R]} = \frac{[L]}{K_d},
\]
we obtain
\[
    P_{\mathrm{bound}} = \frac{[L]/K_d}{1+[L]/K_d}.
\]
Thus
\[
    \boxed{P_{\mathrm{bound}}([L]) = \frac{[L]}{K_d+[L]}.}
\]
This expression has several useful limits.

When there is no ligand,
\[
    [L]=0 \quad \Rightarrow \quad P_{\mathrm{bound}}=0.
\]

When the ligand concentration is very large compared with $K_d$,
\[
    [L]\gg K_d \quad \Rightarrow \quad P_{\mathrm{bound}}\rightarrow1.
\]

When
\[
    [L]=K_d,
\]
we have
\[
    P_{\mathrm{bound}} = \frac{1}{2}.
\]
Thus $K_d$ is the ligand concentration at which half of the receptors are occupied in this simple binding model.
A smaller $K_d$ corresponds to stronger binding, since less ligand is required to reach half occupancy.
% 
% 
% ------------------------------------------------------------
\medskip
\\
\textbf{Step 4: time evolution}
% ------------------------------------------------------------
\\
To solve explicitly for the time evolution, suppose that the free ligand concentration $[L]$ remains approximately constant.
The total receptor concentration is conserved,
\[
    [R]_{\mathrm{tot}} = [R]+[LR].
\]
Therefore,
\[
    [R] = [R]_{\mathrm{tot}}-[LR].
\]
Substituting this into the dynamical equation gives
\[
    \frac{\dd[LR]}{\dd t} = k_{\mathrm{on}}[L]\left([R]_{\mathrm{tot}}-[LR]\right)-k_{\mathrm{off}}[LR].
\]
Collecting terms gives
\[
    \frac{\dd[LR]}{\dd t} = k_{\mathrm{on}}[L][R]_{\mathrm{tot}}-\left(k_{\mathrm{on}}[L]+k_{\mathrm{off}}\right)[LR].
\]
The equilibrium concentration is
\[
    [LR]_{\mathrm{eq}} = [R]_{\mathrm{tot}}\frac{[L]}{K_d+[L]}.
\]
The full time-dependent solution is
\[
    [LR](t) = [LR]_{\mathrm{eq}} + \left([LR](0)-[LR]_{\mathrm{eq}}\right)e^{-\left(k_{\mathrm{on}}[L]+k_{\mathrm{off}}\right)t}.
\]
The system therefore relaxes exponentially toward equilibrium with characteristic timescale
\[
    \tau = \frac{1}{k_{\mathrm{on}}[L]+k_{\mathrm{off}}}.
\]
% 
% 
% ------------------------------------------------------------
\medskip
\\
\textbf{Assumptions of the simple ligand--receptor model}
% ------------------------------------------------------------
\\
The model assumes that:
\begin{itemize}
    \item the system is well mixed, so spatial variation in ligand and receptor concentrations can be ignored;
    \item each receptor has a single independent ligand-binding site;
    \item binding follows the law of mass action;
    \item the rate constants $k_{\mathrm{on}}$ and $k_{\mathrm{off}}$ are constant;
    \item receptors and ligands are chemically identical within their respective populations;
    \item there is no cooperativity between receptors;
    \item there are no additional receptor states or competing ligands;
    \item for the explicit time-dependent solution above, the free ligand concentration $[L]$ is treated as constant.
\end{itemize}


% ============================================================
\subsection{Phytoplankton uptake and resource competition}
% ============================================================

A biological example of constructing dynamical rules from rates is the uptake of limiting resources by phytoplankton.
David Tilman studied competition between freshwater phytoplankton species under different supplies of limiting resources.
The experiments compared predictions from quantitative resource-competition models with the observed abundances of competing algal species.

\begin{quote}
\href{https://hahana.soest.hawaii.edu/cmoreserver/summercourse/2015/documents/Follows_05-30/Tilman-1977.pdf}{
    D. Tilman (1977),
    \textit{Resource Competition between Plankton Algae: An Experimental and Theoretical Approach},
    Ecology \textbf{58}, 338--348.
    }
\end{quote}


% ============================================================
\subsection{Theory of the histogram}
% ============================================================

The theory of the mean follows only the average value of a variable.
In many biological systems, however, the full distribution contains information that the mean alone cannot capture.
Instead of asking only how the mean changes with time, we can ask how probability moves between different possible values or states.

% ------------------------------------------------------------
\subsubsection*{Example 1: number of mRNA molecules in a fly embryo}
% ------------------------------------------------------------

Suppose we want to describe the spatial distribution of mRNA molecules along a fly embryo.
We can discretize the embryo into small spatial boxes.
Each box then contains some number of mRNA molecules.
The collection of counts across boxes forms a spatial distribution rather than a single mean value.


% ------------------------------------------------------------
\subsubsection*{Example 2: age distributions}
% ------------------------------------------------------------

Age is intrinsically continuous, but it can be discretized into bins such as years.

For example,

\[
    [0,1), \qquad [1,2), \qquad [2,3), \qquad \ldots
\]

We can then count how many individuals occupy each age bin.

The resulting histogram is a discretized representation of the underlying age distribution.


% ============================================================
\subsubsection{Theory of the histogram protocol}
% ============================================================

A corresponding protocol for describing an entire probability distribution is:

\begin{enumerate}
    \item Consider the population of interest as a probability distribution over a set of possible states or bins.
    \item Enumerate all the ways in which individuals or probability can move between those bins.
    \item Write an update rule describing the flow of probability into and out of each bin during a time step $\Delta t$.
    \item Given an initial distribution $p(m,t_0)$, solve for the distribution $p(m,t)$ at later times.
\end{enumerate}

Here $m$ labels the state or bin of interest.

For example, $m$ could denote the number of mRNA molecules in a cell.

Instead of following only the mean
\[
    \langle m\rangle,
\]
we now follow the full probability distribution
\[
    p(m,t).
\]
The normalization condition is
\[
    \sum_m p(m,t)=1.
\]
This shift from the dynamics of a mean to the dynamics of a distribution is important whenever fluctuations are themselves biologically relevant.

% ------------------------------------------------------------
\subsubsection{Example: the dynamics of an mRNA histogram}
% ------------------------------------------------------------

Consider a population of identical cells.
Instead of describing only the mean number of mRNA molecules per cell, we describe the entire distribution of mRNA copy numbers.

Let
\[
    p(m,t)
\]
denote the probability that a cell contains exactly $m$ mRNA molecules at time $t$.
Equivalently, if we measure a large population of cells, $p(m,t)$ is the fraction of cells occupying the histogram bin corresponding to $m$ mRNAs.

The distribution is normalized,
\[
    \sum_{m=0}^{\infty}p(m,t)=1.
\]
% % % 
% % % 
% ------------------------------------------------------------
\medskip
\\
\textbf{Step 1: identify the states}
% ------------------------------------------------------------

The possible states are the discrete mRNA copy numbers
\[
    m=0,1,2,3,\ldots
\]
Each value of $m$ corresponds to one bin of the histogram.
% % 
% % 
% ------------------------------------------------------------
\medskip
\\
\textbf{Step 2: enumerate the ways of moving between bins}
% ------------------------------------------------------------

Assume that mRNA is produced at a constant rate $r$.

A transcription event increases the copy number by one,

\[
    m \longrightarrow m+1.
\]

Assume that each individual mRNA molecule is degraded independently at rate $\gamma$.

If there are $m$ molecules in the cell, the total degradation rate is therefore

\[
    \gamma m.
\]

A degradation event decreases the copy number by one,

\[
    m \longrightarrow m-1.
\]

The transitions between neighboring histogram bins are therefore

\[
    \cdots \rightleftharpoons m-1 \rightleftharpoons m \rightleftharpoons m+1 \rightleftharpoons \cdots
\]

with transcription moving probability to the right and degradation moving probability to the left.
% 
% 
% ------------------------------------------------------------
\medskip
\\
\textbf{Step 3: write the probability flow for one bin}
% ------------------------------------------------------------

To determine how $p(m,t)$ changes, enumerate all the ways probability can flow into and out of the bin $m$ during a short time interval $\Delta t$.
There are two ways to enter the state $m$.
First, a cell containing $m-1$ molecules can produce one new mRNA.
The corresponding probability flow is
\[
    r\Delta t\,p(m-1,t).
\]
Second, a cell containing $m+1$ molecules can lose one mRNA through degradation.
Since there are $m+1$ molecules available to degrade, the total degradation rate in this state is $\gamma(m+1)$.
The corresponding probability flow is
\[
    \gamma(m+1)\Delta t\,p(m+1,t).
\]

There are also two ways to leave the state $m$.
A transcription event moves the system from $m$ to $m+1$, giving an outgoing contribution
\[
    r\Delta t\,p(m,t).
\]
A degradation event moves the system from $m$ to $m-1$, giving an outgoing contribution
\[
    \gamma m\Delta t\,p(m,t).
\]
Putting all of these contributions together gives the update rule
\[
    p(m,t+\Delta t)
    \\
    = p(m,t)
    \\
    + r\Delta t\,p(m-1,t)
    \\
    + \gamma(m+1)\Delta t\,p(m+1,t)
    \\
    - r\Delta t\,p(m,t)
    \\
    - \gamma m\Delta t\,p(m,t).
\]
The structure is
\[
    \text{new probability}
    \\
    = \text{old probability}
    \\
    + \text{probability flowing in}
    \\
    - \text{probability flowing out}.
\]
% 
% 
% ------------------------------------------------------------
\medskip
\\
\textbf{Step 4: take the continuous-time limit}
% ------------------------------------------------------------

Subtracting $p(m,t)$ from both sides and dividing by $\Delta t$ gives

\[
    \frac{p(m,t+\Delta t)-p(m,t)}{\Delta t}
    \\
    = rp(m-1,t)
    \\
    + \gamma(m+1)p(m+1,t)
    \\
    - rp(m,t)
    \\
    - \gamma m p(m,t).
\]

Taking the limit $\Delta t\rightarrow0$ gives

\[
    \boxed{
    \frac{\partial p(m,t)}{\partial t}
    = rp(m-1,t)
    + \gamma(m+1)p(m+1,t)
    - rp(m,t)
    - \gamma m p(m,t)
    }.
\]

This is the chemical master equation for constitutive mRNA production and degradation.
It describes the dynamics of the entire histogram rather than only the mean mRNA number.
% 
% 
% ------------------------------------------------------------
\medskip
\\
\textbf{Reading the master equation as probability flow}
% ------------------------------------------------------------

Each term has a direct interpretation:
\begin{enumerate}
    \item $rp(m-1,t)$ : probability flowing into bin $m$ from bin $m-1$ through transcription.
    \item $\gamma(m+1)p(m+1,t)$ : probability flowing into bin $m$ from bin $m+1$ through degradation.
    \item $-rp(m,t)$ : probability leaving bin $m$ through transcription.
    \item $-\gamma m p(m,t)$ : probability leaving bin $m$ through degradation.
\end{enumerate}
Therefore the equation is simply a bookkeeping equation for probability.
% 
% 
% ------------------------------------------------------------
\medskip
\\
\textbf{Boundary at zero mRNAs}
% ------------------------------------------------------------

The copy number cannot become negative.
For $m=0$, there is therefore no probability flow into the state from $m=-1$.
The equation for the zero-mRNA bin is
\[
    \frac{\partial p(0,t)}{\partial t}
    = \gamma p(1,t) - rp(0,t).
\]
The first term represents cells entering the $m=0$ state when their last mRNA is degraded.
The second term represents cells leaving the $m=0$ state when a new mRNA is produced.
% 
% 
% ------------------------------------------------------------
\medskip
\\
\textbf{Steady-state histogram}
% ------------------------------------------------------------

At steady state, the probability of occupying each bin no longer changes with time,
\[
    \frac{\partial p(m,t)}{\partial t}=0.
\]
For this simple birth--death process, the steady-state distribution is a Poisson distribution,
\[
    \boxed{
    p_{\mathrm{ss}}(m)
    = \frac{\lambda^m e^{-\lambda}}{m!}
    }.
\]
The parameter $\lambda$ is
\[
    \lambda=\frac{r}{\gamma}.
\]
Therefore,
\[
    \boxed{
    p_{\mathrm{ss}}(m)
    = \frac{1}{m!}
    \left(\frac{r}{\gamma}\right)^m
    e^{-r/\gamma}
    }.
\]
The mean of this distribution is
\[
    \langle m\rangle_{\mathrm{ss}}=\frac{r}{\gamma}.
\]
Its variance is also
\[
    \operatorname{Var}(m)_{\mathrm{ss}}=\frac{r}{\gamma}.
\]
% 
% 
% ------------------------------------------------------------
\medskip
\\
\textbf{Connection to the theory of the mean}
% ------------------------------------------------------------

The histogram description contains the mean description as one of its consequences.
The mean mRNA number is defined from the probability distribution as
\[
    \langle m(t)\rangle = \sum_{m=0}^{\infty}m\,p(m,t).
\]
Using the master equation, one finds
\[
    \boxed{
    \frac{\dd\langle m\rangle}{\dd t}
    = r-\gamma\langle m\rangle
    }.
\]
This is exactly the deterministic equation obtained using the theory-of-the-mean protocol.
At steady state,
\[
    \langle m\rangle_{\mathrm{ss}}=\frac{r}{\gamma}.
\]

The two approaches therefore answer different levels of the same question.
The theory of the mean tells us how the average copy number evolves.
The theory of the histogram tells us how the probability of every possible copy number evolves.
The histogram description consequently contains information about fluctuations, variance, and the shape of the population distribution that is invisible if we follow only the mean.
% 
% 
% ------------------------------------------------------------
\medskip
\\
\textbf{Assumptions}
% ------------------------------------------------------------

The simple mRNA birth-death model assumes that:

\begin{itemize}
    \item transcription occurs at a constant rate $r$ independent of the current mRNA number;
    \item each mRNA molecule degrades independently at the same rate $\gamma$;
    \item transcription changes the copy number by exactly $+1$;
    \item degradation changes the copy number by exactly $-1$;
    \item the system is Markovian, so transition probabilities depend only on the current state;
    \item $\Delta t$ is sufficiently small that the probability of two or more reaction events during the same time interval is of order $\Delta t^2$ and can be neglected;
    \item there is no transcriptional bursting, promoter switching, feedback, or other regulation;
    \item cells are treated as identical apart from the stochastic differences generated by transcription and degradation.
\end{itemize}


% ============================================================
\subsection{Biological turnover}
% ============================================================

Biological systems are continually renewed through the production, loss, and replacement of their components.

Turnover can therefore be viewed using the same general framework: identify the states of interest, enumerate the processes that move material between those states, assign rates, and construct the corresponding dynamical equations.

Reference on biological turnover:
\url{https://www.sciencedirect.com/science/article/pii/S0092867424004604}

Reference on red blood cell turnover:
\url{https://www.sciencedirect.com/science/article/pii/S0021925817352018}
